\documentclass[11pt]{article}
\usepackage[margin=1in]{geometry}
\usepackage{amsmath,amssymb,amsthm,mathrsfs}
\usepackage{graphicx}
\usepackage[hidelinks]{hyperref}

\newtheorem*{theorem}{Theorem}
\newtheorem*{lemma}{Lemma}
\newtheorem*{corollary}{Corollary}

\title{A Complemented-Subspace Reduction for a Rosenthal--Odell Quotient Question}
\author{Run \texttt{fa\_banach\_001}, lane 8}
\date{June 14, 2026}

\begin{document}
\maketitle

\section*{Status}

\textbf{partial\_result (complemented subspaces of quotients of strict \(\ell_1\)-preduals contain \(c_0\), plus a general \((WU)\) trichotomy and \(c_0\)-saturation criterion).}

This packet is the single durable packet for the current Rosenthal--Odell target.  Its promoted mathematical claim is the partial theorem below.  The later sections record additional constraints on any remaining counterexample, including a general \((WU)\)-based trichotomy.  These reductions are not a full solution.  A broader literature-implied answer for the isomorphic \(\ell_1\)-predual interpretation is recorded later as context, not as the promoted claim.

\section*{Source Question}

The source paper is Gasparis, \emph{New examples of \(c_0\)-saturated Banach spaces}, arXiv:0809.1808 \cite{Gasparis}.  On page 2, Gasparis records Rosenthal's question asking whether every quotient of an \(\ell_1\)-predual is \(c_0\)-saturated, and Odell's special case for quotients of \(C(K)\) with \(K\) countable compact.

\begin{center}
\includegraphics[width=\textwidth]{figures/open_problem_crop.png}
\end{center}

\section*{Partial Result}

\begin{theorem}
Let \(E\) be a real Banach space such that \(E^*\) is isometric to \(\ell_1\), and let \(X\) be a quotient of \(E\).  Then every complemented infinite-dimensional subspace of \(X\) contains an isomorphic copy of \(c_0\).

Consequently, if \(K\) is a countable compact metric space and \(X\) is a quotient of \(C(K)\), every complemented infinite-dimensional subspace of \(X\) contains \(c_0\).
\end{theorem}

\section*{Proof Intuition}

The full Rosenthal--Odell problem is hard because a bad subspace of a quotient need not itself be a quotient.  The complemented case removes exactly that difficulty.  If \(Y\) is complemented in \(X\), the projection \(P:X\to Y\) turns \(Y\) into the range of an operator on the whole quotient.  Quotients of real \(L^1\)-preduals have Pelczynski's property \((V)\), so if \(Y\) contains no \(c_0\), the projection must be weakly compact.  That makes \(Y\) reflexive.  But in the strict \(\ell_1\)-predual/countable \(C(K)\) setting, a complemented subspace of a quotient is itself a quotient, so its dual embeds into \(\ell_1\); an infinite-dimensional reflexive space cannot have dual embedded in \(\ell_1\).  Thus the complemented bad subspace cannot exist.

\section*{Proof}

We use four standard facts.  First, by Johnson--Zippin's theorem, in the form recalled for example in \cite[Theorem 4.5]{Krulisova}, every real \(L^1\)-predual has Pelczynski's property \((V)\).  Second, property \((V)\) passes to quotients: if \(Q:E\to X\) is onto and \(T:X\to Z\) is unconditionally converging, then \(TQ:E\to Z\) is unconditionally converging, hence weakly compact; by the open mapping theorem this implies that \(T\) is weakly compact.  Third, an operator is unconditionally converging if and only if it fixes no copy of \(c_0\).  Fourth, every subspace of \(\ell_1\) has the Schur property, so it contains no infinite-dimensional reflexive subspace.

Let \(Q:E\to X\) be a quotient map and let \(Y\subset X\) be complemented.  Let \(P:X\to Y\) be a bounded projection.  Suppose, toward a contradiction, that \(Y\) is infinite-dimensional and contains no copy of \(c_0\).

The projection \(P\) fixes no copy of \(c_0\).  Indeed, if there were a subspace \(Z\subset X\), \(Z\simeq c_0\), such that \(P|_Z\) were an isomorphic embedding, then \(PZ\subset Y\) would be isomorphic to \(c_0\), contrary to the assumption on \(Y\).  Hence \(P\) is unconditionally converging.  Since \(X\) is a quotient of the real \(L^1\)-predual \(E\), the quotient-stability observation above gives property \((V)\) for \(X\), and therefore \(P\) is weakly compact.

But \(P|_Y\) is the identity on \(Y\).  Thus \(B_Y\subset P(B_X)\), up to the harmless inclusion \(B_Y\subset B_X\), and \(B_Y\) is relatively weakly compact.  Hence \(Y\) is reflexive.

Now \(P Q:E\to Y\) is onto.  Its adjoint \((PQ)^*:Y^*\to E^*\) is an isomorphic embedding.  Since \(E^*\) is isometric to \(\ell_1\), \(Y^*\) is isomorphic to a reflexive subspace of \(\ell_1\).  A reflexive Schur space is finite-dimensional: its closed unit ball is weakly compact and, by the Schur property, weak compactness is norm compactness.  Hence \(Y^*\), and therefore \(Y\), is finite-dimensional, contradicting the choice of \(Y\).

This proves that every complemented infinite-dimensional subspace of \(X\) contains \(c_0\).  If \(K\) is countable compact metric, then \(C(K)^*\) is isometric to \(\ell_1(K)\), so the same conclusion applies to every quotient of \(C(K)\).

\section*{Literature-Implied Broad Answer}

There is a broader reading of Rosenthal's phrase ``\(\ell_1\)-predual'' under which one only requires \(X^*\) to be isomorphic to \(\ell_1\).  On that broader, isomorphic reading, the quotient question has a negative answer by direct identification with later Bourgain--Delbaen/Argyros--Motakis constructions.

Argyros--Motakis \cite{AM} construct a hereditarily indecomposable \(\mathscr L_\infty\)-space \(\mathfrak X_{\mathrm{nr}}\) which contains no \(c_0\), no \(\ell_1\), and no reflexive subspace.  In their quotient section they state that there are infinite-dimensional spaces \(X_1,X_2,X_3\) such that \(X_2\) is a quotient of \(X_1\), \(X_3\) is a quotient of \(X_2\), \(X_1\) and \(X_3\) are reflexively saturated, \(X_2\) contains no reflexive subspace, and all three spaces are \(\ell_1\)-preduals with the scalar-plus-compact property.  The surrounding discussion identifies the middle space with the non-\(c_0\) \(\mathscr L_\infty\) construction.

Thus, if ``\(\ell_1\)-predual'' is interpreted isomorphically, \(X_1\twoheadrightarrow X_2\) is a quotient of an \(\ell_1\)-predual which is not \(c_0\)-saturated, since \(X_2\) itself contains no copy of \(c_0\).  This is a literature-implied answer and is not claimed here as a new proof.  It also does not settle Odell's special problem for quotients of \(C(K)\), \(K\) countable compact.

\section*{Final-Push Reduction for the Countable \(C(K)\) Case}

The following elementary observation is useful when trying to attack the uncomplemented part of Odell's problem.

\begin{lemma}
Let \(K\) be countable compact metric, let \(X\) be a quotient of \(C(K)\), and let \(Y\subset X\) be a subspace.  Then \(Y^*\) is separable.  Consequently \(Y\) contains no copy of \(\ell_1\).
\end{lemma}

\begin{proof}
Since \(X\) is a quotient of \(C(K)\), \(X^*\) embeds isometrically into \(C(K)^*\cong \ell_1(K)\).  In particular \(X^*\) is separable.  The restriction map \(X^*\to Y^*\) is onto by Hahn--Banach, so \(Y^*\) is separable as well.  If \(Y\) contained \(\ell_1\), then restriction from \(Y^*\) onto the dual of that copy would give a quotient map \(Y^*\to \ell_\infty\).  This is impossible because quotients of separable spaces are separable while \(\ell_\infty\) is not.
\end{proof}

For the first hard ordinal level \(C(\omega^\omega)\), Odell records a further obstruction due to Schlumprecht \cite{Odell}: if \(X\) is a quotient of \(C(\omega^\omega)\), then the closed linear span of any normalized weakly null sequence in \(X\) whose spreading model is equivalent to the \(\ell_1\) basis contains \(c_0\).  Therefore any \(c_0\)-free subspace \(Y\) of a quotient of \(C(\omega^\omega)\) must be very special: it has separable dual, contains no \(\ell_1\), and has no normalized weakly null sequence with \(\ell_1\) as a spreading model.

This reduction explains why the remaining case resembles the modern Bourgain--Delbaen/Argyros--Haydon low-Szlenk pathology rather than a simple \(\ell_1\) obstruction.  Alspach's analysis of the Bourgain--Delbaen space \cite{AlspachBD} shows that there are \(\mathscr L_\infty\)-spaces with separable dual, no \(c_0\), and small Szlenk index; such examples stress-test any proposed proof, even though they are not known here to be quotients of countable ordinal \(C\)-spaces.  Conversely, Alspach's earlier note \cite{AlspachPredual} emphasizes that quotients of \(C(\alpha)\), \(\alpha<\omega_1\), were already subtle in the strict \(\ell_1\)-predual setting.

\section*{Extra Push: General \((WU)\) Trichotomy}

Odell \cite{Odell} says that a Banach space has property \((WU)\) when every normalized weakly null sequence has an unconditional subsequence.  His Theorem A proves that every quotient of a space with a shrinking unconditional finite-dimensional decomposition has \((WU)\).  The same paper also points out that \(C(\omega^\omega)\) has \((WU)\), but asks whether quotients preserve \((WU)\) in this setting.

The preceding countable \(C(K)\) reduction is a special case of the following general trichotomy.  The point is not a new construction, but a useful packaging of Rosenthal's theorem, Odell's \((WU)\) hypothesis, and James's unconditional-basis criteria.

\begin{theorem}[General \((WU)\) trichotomy]
Let \(X\) be an infinite-dimensional Banach space with property \((WU)\).  Then every infinite-dimensional subspace \(Y\subset X\) contains a copy of \(\ell_1\), a copy of \(c_0\), or an infinite-dimensional reflexive subspace.

In particular, if \(X\) has \((WU)\) and contains neither \(\ell_1\) nor \(c_0\), then \(X\) is reflexively saturated.
\end{theorem}

\begin{proof}
Property \((WU)\) passes to subspaces: if \(Y\subset X\) and a sequence is weakly null in \(Y\), then it is weakly null in \(X\).  Let \(Y\subset X\) be infinite-dimensional.  If \(Y\) contains \(\ell_1\) or \(c_0\), there is nothing to prove.  Assume that \(Y\) contains neither \(\ell_1\) nor \(c_0\).

If \(Y\) is reflexive, again there is nothing to prove.  Suppose that \(Y\) is not reflexive.  By Eberlein--Smulian, choose a bounded sequence in \(Y\) with no weakly convergent subsequence.  Since \(Y\) contains no \(\ell_1\), Rosenthal's \(\ell_1\)-theorem \cite{RosenthalL1} gives a weakly Cauchy subsequence \((y_n)\).  This subsequence is not weakly convergent, by the choice of the original sequence, and hence it is not norm Cauchy.

Thus there are indices \(n_1<m_1<n_2<m_2<\cdots\) and a number \(\delta>0\) such that
\[
  \|y_{m_j}-y_{n_j}\|\ge \delta\qquad (j\in\mathbb N).
\]
Set
\[
  z_j=\frac{y_{m_j}-y_{n_j}}{\|y_{m_j}-y_{n_j}\|}.
\]
For every \(f\in Y^*\), the scalar sequence \(f(y_n)\) is Cauchy, so \(f(z_j)\to 0\).  Hence \((z_j)\) is normalized and weakly null in \(Y\).  By \((WU)\), some subsequence \((z_{j_k})\) is unconditional basic.

Let \(Z=[z_{j_k}]\).  Since \(Z\subset Y\), the space \(Z\) contains neither \(\ell_1\) nor \(c_0\).  James's characterizations for unconditional bases now apply \cite{JamesBases}: an unconditional basis is shrinking if its closed span contains no \(\ell_1\), and boundedly complete if its closed span contains no \(c_0\).  Therefore \((z_{j_k})\) is both shrinking and boundedly complete, so \(Z\) is reflexive.  This gives the required infinite-dimensional reflexive subspace of \(Y\).
\end{proof}

\begin{corollary}[No-\(\ell_1\) \((WU)\) criterion for \(c_0\)-saturation]
Let \(X\) have property \((WU)\) and contain no copy of \(\ell_1\).  Then every infinite-dimensional subspace of \(X\) contains \(c_0\) or an infinite-dimensional reflexive subspace.  Consequently,
\[
  X\text{ is }c_0\text{-saturated}
  \quad\Longleftrightarrow\quad
  X\text{ contains no infinite-dimensional reflexive subspace.}
\]
\end{corollary}

\begin{proof}
The first assertion is the trichotomy with the \(\ell_1\) alternative removed.  If \(X\) contains an infinite-dimensional reflexive subspace, then \(X\) is not \(c_0\)-saturated, since subspaces of reflexive spaces are reflexive and \(c_0\) is not reflexive.  Conversely, if \(X\) is not \(c_0\)-saturated, choose an infinite-dimensional subspace \(Y\subset X\) containing no \(c_0\).  The first assertion gives an infinite-dimensional reflexive subspace of \(Y\), and hence of \(X\).
\end{proof}

\begin{corollary}[Quotients of shrinking unconditional FDD spaces]
Let \(X\) be a quotient of a Banach space with a shrinking unconditional finite-dimensional decomposition.  Then every infinite-dimensional subspace of \(X\) contains \(c_0\) or an infinite-dimensional reflexive subspace.  In particular, \(X\) is \(c_0\)-saturated if and only if \(X\) contains no infinite-dimensional reflexive subspace.
\end{corollary}

\begin{proof}
By Odell's theorem \cite{Odell}, \(X\) has property \((WU)\).  The dual of the source space is separable, so \(X^*\), which embeds into that dual, is separable.  Therefore \(X\) contains no copy of \(\ell_1\).  Apply the preceding corollary.
\end{proof}

For the countable compact \(C(K)\) target, the same criterion applies whenever the quotient \(X\) itself is known to have \((WU)\): the earlier lemma gives \(X^*\) separable, hence no \(\ell_1\), and the no-\(\ell_1\) \((WU)\) criterion says that \(X\) is \(c_0\)-saturated exactly when \(X\) contains no infinite-dimensional reflexive subspace.  More locally, any \(c_0\)-free subspace \(Y\) of a quotient of countable \(C(K)\) which has \((WU)\) must contain an infinite-dimensional reflexive subspace.

For quotients of \(C(\omega^\omega)\), this combines with Schlumprecht's observation as follows.  A hypothetical \(c_0\)-free subspace with no infinite-dimensional reflexive subspace must fail \((WU)\): it contains a normalized weakly null sequence with no unconditional subsequence.  At the same time, none of its normalized weakly null sequences can have an \(\ell_1\) spreading model, because such a sequence would generate a closed span containing \(c_0\).  Thus any no-reflexive counterexample must be a very specific Maurey--Rosenthal-type obstruction inside a quotient of \(C(\omega^\omega)\), not merely a Bourgain--Delbaen-style no-\(c_0\) space with separable dual.

\section*{Limitations}

This is not a full solution of the Rosenthal--Odell question.  The open problem asks about every infinite-dimensional subspace of the quotient.  The argument above uses the projection onto \(Y\), so it only rules out complemented \(c_0\)-free subspaces.  Any counterexample to the strict \(\ell_1\)-predual or countable compact \(C(K)\) version must therefore contain its \(c_0\)-free subspace in an uncomplemented way.

A naive Cantor--Bendixson induction would quotient by the first derivative kernel \(C_0(K\setminus K')\cong c_0\).  The induction would close if the image of this kernel in every quotient were closed or at least \(c_0\)-saturated.  This is exactly the dangerous point: bounded operators from \(c_0\) may have nonclosed range with arbitrary separable closure, so the closure of the kernel image can in principle hide a \(c_0\)-free low-Szlenk subspace.  Any full proof must rule out that behavior using the specific ordinal \(C(K)\) structure, not just the abstract exact sequence.

Another tempting shortcut also fails: although \(Y^*\) is a quotient of a subspace of \(\ell_1\), it need not have the Schur property.  Quotients of \(\ell_1\) can be arbitrary separable Banach spaces.  Thus the final-push reduction cannot rule out reflexive or James-type subspaces merely from the dual embedding into \(\ell_1\); it only rules out \(\ell_1\) itself and, in the \(C(\omega^\omega)\) case, weakly null \(\ell_1\)-spreading-model routes.

The extra-push \((WU)\) trichotomy above also does not solve the problem.  It shows that, under \((WU)\) and no \(\ell_1\), failure of \(c_0\)-saturation is equivalent to the presence of an infinite-dimensional reflexive subspace.  Odell explicitly lists quotient preservation of \((WU)\) as open in the relevant \(C(\omega^\omega)\) neighborhood, and the existence of reflexive subspaces in quotients of countable \(C(K)\) spaces is not ruled out by the present argument.

\section*{Verification Notes}

The proof is qualitative and uses no computation.  The key review points are the real-scalar convention in Johnson--Zippin's property \((V)\) theorem and the standard equivalence between unconditionally converging operators and operators fixing no copy of \(c_0\).  The countable compact \(C(K)\) specialization is strict: \(C(K)^*\cong \ell_1(K)\) isometrically.

For the literature-implied broad answer, review should check the convention for ``\(\ell_1\)-predual'' and the identification of the middle quotient \(X_2\) in \cite{AM} with the no-\(c_0\) Argyros--Motakis space.  That contextual answer should not be used to mark the countable compact \(C(K)\) problem solved.

For the final-push reduction, review should check two points.  First, the separable-dual obstruction uses only \(X^*\subset \ell_1(K)\) and Hahn--Banach restriction, and deliberately does not claim that \(Y^*\) has the Schur property.  Second, Schlumprecht's observation as recorded in \cite{Odell} is used here only for quotients of \(C(\omega^\omega)\) and only to rule out \(\ell_1\)-spreading-model weakly null sequences inside a hypothetical \(c_0\)-free subspace.

For the extra-push \((WU)\) trichotomy, review should check the inheritance of \((WU)\) by subspaces, the use of Rosenthal's \(\ell_1\)-theorem to obtain a nontrivial weakly Cauchy sequence, and James's shrinking/boundedly-complete characterizations for unconditional bases.  The criterion assumes \((WU)\) for the bad subspace or quotient; it does not assert that arbitrary quotients of \(C(\omega^\omega)\) have \((WU)\).  A bounded local-index and web search found Odell's source theorem and related problem context, but not an exact pre-existing statement of the packaged trichotomy; since the proof is a direct combination of standard theorems, novelty should be treated as a useful derived reduction rather than as a claimed independent classical theorem.

\begin{thebibliography}{9}

\bibitem{Gasparis}
I. Gasparis,
\emph{New examples of \(c_0\)-saturated Banach spaces},
arXiv:0809.1808, 2008.

\bibitem{Krulisova}
H. Krulisova,
\emph{Quantification of Pelczynski's property (V)},
arXiv:1509.06610, 2015.

\bibitem{AM}
S. A. Argyros and P. Motakis,
\emph{The scalar-plus-compact property in spaces without reflexive subspaces},
arXiv:1608.01962, 2016.

\bibitem{Odell}
E. Odell,
\emph{On quotients of Banach spaces having shrinking unconditional bases},
Illinois J. Math. \textbf{36} (1992), 681--695.

\bibitem{RosenthalL1}
H. P. Rosenthal,
\emph{A characterization of Banach spaces containing \(\ell_1\)},
Proc. Nat. Acad. Sci. U.S.A. \textbf{71} (1974), 2411--2413.

\bibitem{JamesBases}
R. C. James,
\emph{Bases and reflexivity of Banach spaces},
Ann. of Math. (2) \textbf{52} (1950), 518--527.

\bibitem{AlspachPredual}
D. E. Alspach,
\emph{A \(\ell_1\)-predual which is not isometric to a quotient of \(C(\alpha)\)},
arXiv:math/9204215, 1992.

\bibitem{AlspachBD}
D. E. Alspach,
\emph{The dual of the Bourgain--Delbaen space},
arXiv:math/9805081, 1998.

\end{thebibliography}

\end{document}
